\chapter{A Concrete Naming Certificate for $\CompactMetricLocatedNetSelectionUp$}
\label{ch:concrete-instances-compact-metric-located-net-selection-namecert}
\origin{human}

Following Bishop and Bridges constructive compact metric analysis, the $\CompactMetricLocatedNetSelectionUp$ packet records the finite located-net selection surface consumed by L10 regular-sequence and Real readback routes. It sits on the compact-metric source of \autoref{ch:concrete-instances-compactmetric-namecert}, the dyadic tolerance ledger of \autoref{ch:concrete-instances-dyadicratcore-namecert}, the finite observation windows of \autoref{ch:concrete-instances-streamname-namecert}, the regular rational readback of \autoref{ch:concrete-instances-regseqrat-namecert}, and the terminal real-name seal of \autoref{ch:concrete-instances-real-namecert}. The packet selects only finite located-net rows already displayed by the carrier; it does not assert open-cover compactness, selected subcovers, quotient equality, countable choice, or ambient $\RealUp$ completeness.

\begin{definition}[Compact metric located-net selection carrier]
\label{def:compact-metric-located-net-selection-carrier}
A $\CompactMetricLocatedNetSelectionUp$ carrier is a finite $\BHist$ packet
$$
L=(K,D,W,R,E,H,C,P,N)
$$
with the following displayed rows:
\begin{enumerate}
\item $K$ is the $\CompactMetricUp$ source row naming the compact metric packet whose finite net may be inspected.
\item $D$ is the $\DyadicRatCoreUp$ tolerance row against which located distances and selected net radii are read.
\item $W$ is the $\StreamNameUp$ finite-window row opened by the tolerance request.
\item $R$ is the $\RegSeqRatUp$ readback row that transports the located net window into the regular rational route.
\item $E$ is the terminal $\RealUp$ seal row reached only after $K,D,W,R$ have been displayed.
\item $H$ records componentwise $\hsame$ transport for compact-metric, dyadic, window, readback, seal, replay, provenance, and local naming rows.
\item $C$ records displayed $\Cont$ replay from a located-net request through $K,D,W,R$ and into the terminal seal $E$.
\item $P$ is the $\Pkg$ provenance over $\CompactMetricLocatedNetSelectionUp$, $\CompactMetricUp$, $\DyadicRatCoreUp$, $\StreamNameUp$, $\RegSeqRatUp$, $\RealUp$, $\BHist$, $\hsame$, $\Cont$, and $\NameCert$.
\item $N$ is the local $\NameCert_{\CompactMetricLocatedNetSelectionUp}$ row.
\end{enumerate}
The carrier has no coordinate for open-cover compactness, selected subcovers, quotient real equality, countable choice, an infinite choice sequence of net points, or ambient Real completeness.
\end{definition}

\begin{theorem}[Compact metric located-net selection obligations]
\label{thm:compact-metric-located-net-selection-obligations}
For every accepted $\CompactMetricLocatedNetSelectionUp$ carrier $L=(K,D,W,R,E,H,C,P,N)$, the local $\NameCert_{\CompactMetricLocatedNetSelectionUp}$ surface consists exactly of compact-metric source exposure through $K$, dyadic tolerance exposure through $D$, finite StreamName window exposure through $W$, RegSeqRat readback through $R$, terminal Real sealing through $E$, componentwise $\hsame$ transport through $H$, displayed $\Cont$ replay through $C$, $\Pkg$ provenance through $P$, and local naming through $N$.
\end{theorem}
\begin{proof}
Project the carrier. The compact-metric row $K$ supplies the only source of finite net data, while $D,W,R,E$ are the only analytic rows available before a Real-facing read is exposed.

Read the located-net request through the displayed L10 route. The dyadic tolerance row fixes the finite radius request, the StreamName row opens only the needed observation window, and the RegSeqRat row records the regular rational readback for that window.

The structural tail is exhausted by $H,C,P,N$. Componentwise transport, continuation replay, provenance, and local naming preserve the same finite packet, so the NameCert obligations cannot export open-cover compactness, selected subcovers, quotient equality, countable choice, or ambient Real completeness.
\end{proof}

\begin{theorem}[Compact metric located-net selection L10 handoff]
\label{thm:compact-metric-located-net-selection-l10-handoff}
Every L10 regular-sequence or Real readback consumer of an accepted $\CompactMetricLocatedNetSelectionUp$ carrier must first read the finite route
$$
\CompactMetricUp \longmapsto \DyadicRatCoreUp \longmapsto \StreamNameUp \longmapsto \RegSeqRatUp \longmapsto \RealUp .
$$
The handoff exports a located finite-net window over displayed compact-metric and dyadic rows, not open-cover compactness, a selected subcover, quotient real equality, countable choice, or ambient Real completeness.
\end{theorem}
\begin{proof}
Begin with the NameCert obligations. They expose $K,D,W,R,E,H,C,P,N$ before any consumer can treat the located-net selection as available.

The L10 route is ordered by projection from the carrier. The compact metric row is read first, the dyadic tolerance row determines the requested finite radius, the StreamName window supplies the inspected observations, and the RegSeqRat row performs the regular rational readback.

Only after those rows are visible may the terminal Real seal be consumed. Since no coordinate stores open-cover compactness, selected subcovers, quotient equality, countable choice, or ambient Real completeness, the handoff remains a finite located-net readback surface.
\end{proof}

\closureat{CompactMetricLocatedNetSelectionUp}{seedStr}

\begin{closurestatus}{\CompactMetricLocatedNetSelectionUp}
  \constructivestory{A $\CompactMetricLocatedNetSelectionUp$ packet is generated from compact-metric source rows, dyadic tolerance rows, finite StreamName windows, RegSeqRat readback, terminal Real sealing, componentwise hsame transport, displayed Cont replay, Pkg provenance, and local NameCert rows.}
  \theoryclosure{\seedClosure}
  \scopeclosed{\autoref{def:compact-metric-located-net-selection-carrier}, \autoref{thm:compact-metric-located-net-selection-obligations}, and \autoref{thm:compact-metric-located-net-selection-l10-handoff} seed the finite located-net selection surface for L10 RegSeqRat and Real readback consumers.}
  \formalstatus{\unformalizedV}
  \leantarget{BEDC.Derived.CompactMetricLocatedNetSelectionUp}
  \bridgestatus{none}
  \notclaimed{This chapter does not claim open-cover compactness, selected subcovers, quotient real equality, countable choice, an infinite choice sequence of net points, ambient Real completeness, Lean theorem verification, or bridge checking.}
  \upgradepath{Obligation closure requires checked carrier admission, located-net row exposure, dyadic tolerance routing, finite-window readback, terminal Real sealing, transport, replay, provenance, and local naming for BEDC.Derived.CompactMetricLocatedNetSelectionUp.}
\end{closurestatus}
